\section{Complete solver ablation}
\label{app:solver}

Variant \textbf{A} uses only the hand-written rule path, \textbf{B} only
retrieval, and \textbf{C} is the deployed configuration, which tries the rule
path first and falls back to retrieval. If the solver's accuracy came from
copying gold programs out of its retrieval library it would be a memorization
artifact rather than a baseline, so the quantity of interest is $C-A$: what
retrieval adds on top of the rules.

\begin{table}[t]\centering\small
\setlength{\tabcolsep}{4pt}
\caption{Complete solver ablation. \textbf{A} uses only the hand-written rule
path, \textbf{B} only retrieval, \textbf{C} is the deployed configuration
(rules first, retrieval as fallback). Dev accuracy (\%), mean over seeds, at
every budget of each configuration. Retrieval never fires at any budget used
for a reported group.}
\label{tab:supp-ablation}
\begin{tabular}{cr rrr r}
\toprule
cfg & $m$ & A: rules & B: retr. & C: deployed & $C-A$ \\
\midrule
1 & 0 & 72.00 & 0.00 & 72.00 & $+0.00$ \\
 & 1 & 72.00 & 0.00 & 72.00 & $+0.00$ \\
 & 2 & 72.00 & 0.00 & 72.00 & $+0.00$ \\
 & 4 & 72.00 & 0.00 & 72.00 & $+0.00$ \\
 & 8 & 72.00 & 0.00 & 72.00 & $+0.00$ \\
 & 16 & 72.00 & 0.00 & 72.00 & $+0.00$ \\
 & 64 & 72.00 & 0.00 & 72.00 & $+0.00$ \\
 & 256 & 72.00 & 0.00 & 72.00 & $+0.00$ \\
 & 1024 & 72.00 & 0.00 & 72.00 & $+0.00$ \\
 & 2446 & 72.00 & 0.00 & 72.00 & $+0.00$ \\
2 & 0 & 34.50 & 0.00 & 34.50 & $+0.00$ \\
 & 16 & 35.67 & 0.00 & 35.50 & $-0.17$ \\
 & 64 & 35.83 & 0.00 & 35.17 & $-0.67$ \\
 & 128 & 36.00 & 0.00 & 35.33 & $-0.67$ \\
 & 256 & 35.67 & 0.00 & 35.33 & $-0.33$ \\
 & 1024 & 36.00 & 1.00 & 36.33 & $+0.33$ \\
 & 3635 & 36.00 & 2.00 & 38.00 & $+2.00$ \\
\bottomrule
\end{tabular}
\end{table}


\noindent
Variant~B returns $0.00\%$ at every one of Config~1's ten budgets and at every
Config~2 budget up to and including $m{=}256$: it never produces an accepted
program. Across all ten Config~1 budgets and three seeds---thirty
measurements---$A$, $C$ and their difference are constant at $72.00\%$,
$72.00\%$ and $+0.00$ points.

Config~2 is the only place retrieval ever fires, and it does so only at
$m{=}1024$ ($1.00\%$) and at the full pool $m{=}3635$ ($2.00\%$). Both lie far
above the two operating points from which reported groups are drawn
($m_L{=}128$, $m_H{=}256$); at those two the deployed variant differs from the
rule-only path by $-0.67$ and $-0.33$ points, which is one to two items on the
200-item dev subset.

With retrieval never firing at any operating budget, the demonstration budget
has no channel through which to affect the solver at all. That is the mechanism
behind the flat dose curve in the main paper: the solver is not merely
insensitive to $m$ in practice, it is structurally unable to use $m$.
